\section{实证补充结果与检验方法}
\label{app:empirical}
本附录保存正文所用配置的补充结果、风险匹配口径和统计检验。各项计算沿用正文的数据与实施规则，完整逐期持仓保存在仓库 CSV 文件中。
\subsection{换手与成本}
Global RRP 月均买卖合计换手为 \globalMonthlyTurnover。图\ref{fig:turnover}比较共同周度日历下各配置规则的交易需求。交易量以目标权重和交易前漂移权重之差计算，能够反映主动调整及价格变化带来的再平衡需求。

\figmaybe{global_rrp_turnover_comparison.pdf}{五模型月均买卖合计换手}{fig:turnover}

主模型毛年化收益为 \globalGrossReturn，扣除单边 \txCostBps{} bps 费用后，净年化收益为 \globalNetReturn，年化收益差为 \globalCostDrag。年度选择中的换手规则作用于参数筛选，最终交易成本由实际持仓路径决定。图\ref{fig:cost_drag}报告各模型毛、净收益复利路径的年化差值；后文通过多档费率进一步量化成本敏感性。

\figmaybe{global_rrp_transaction_cost_comparison.pdf}{五模型约定交易成本的年化收益拖累}{fig:cost_drag}
\FloatBarrier

\subsection{尾部损失}
图\ref{fig:cvar}报告95\%日度历史CVaR，Global RRP 的对应值为 \globalCVaR。该指标汇总日损失分布尾部的平均水平，与最大回撤共同描述下行风险。CVaR 用于事后评价，优化中的风险项采用归一化方差。

\figmaybe{global_rrp_cvar_comparison.pdf}{五模型历史日度尾部损失}{fig:cvar}

\subsection{消融与参数冻结的完整结果}
\begin{table}[htbp]\centering
\caption{目标项消融与参数冻结的完整比较}\label{tab:ablation_full}
\fitthesistable{\begin{tabular}{lrrrrr}
\toprule
实验设定 & 净年化收益 & 年化波动 & 夏普 & 最大回撤 & 月均换手\\\midrule
Global RRP & 5.98\% & 3.58\% & 1.641 & -6.80\% & 17.29\%\\
仅风险预算参考 & 5.87\% & 4.76\% & 1.222 & -7.72\% & 9.96\%\\
收益短缺惩罚置零 & 5.38\% & 3.37\% & 1.574 & -7.08\% & 11.36\%\\
方差惩罚置零 & 5.87\% & 4.76\% & 1.222 & -7.72\% & 9.95\%\\
冻结评价首年参数 & 5.96\% & 3.55\% & 1.651 & -6.58\% & 17.67\%\\
\bottomrule
\end{tabular}}
\end{table}
\begin{table}[htbp]\centering
\caption{资本跟踪与风险贡献偏离}\label{tab:diagnostic_mechanism}
\fitthesistable{\begin{tabular}{lrrr}
\toprule
实验设定 & 平均权重距离平方 & 平均最大风险预算偏离 & 平均货币权重\\\midrule
Global RRP & 0.00232 & 6.08\% & 18.02\%\\
仅风险预算参考 & 0.00000 & 0.00\% & 16.23\%\\
收益短缺惩罚置零 & 0.00244 & 6.53\% & 18.14\%\\
方差惩罚置零 & 0.00000 & 0.00\% & 16.23\%\\
冻结评价首年参数 & 0.00244 & 6.34\% & 18.10\%\\
\bottomrule
\end{tabular}}
\end{table}

表\ref{tab:diagnostic_mechanism}将资本权重距离和风险预算偏离分别列示。完整模型的平均最大风险预算偏离为\diagPrimaryRCError{}个百分点，参考组合对应\diagReferenceRCError{}个百分点。计算先在每个决策日取资产风险贡献份额相对等预算的最大绝对差，再对决策日取平均；它衡量预算偏离，并非组合年化波动或收益损失。该指标统一使用主模型当期Ledoit--Wolf矩阵，参考组合接近等预算，最终组合则体现收益与方差调整后的偏离。

\figmaybe{primary_diagnostic_ablation.pdf}{目标项消融的收益与风险预算偏离}{fig:diagnostic_ablation}

冻结评价首年参数的净年化收益为 \diagFrozenReturn，夏普为 \diagFrozenSharpe，月均换手为 \diagFrozenTurnover。该实验把参数更新的作用与凸优化结构本身区分开来。年度更新与冻结规则的差异结合下文配对抽样区间评价，年度候选集合的宽度则用于说明选择规则的统计分辨能力。
\FloatBarrier


\subsection{事前风险匹配的计算与可行性}
对 Equal Weight 和 60/40 Benchmark，记其当期目标为 $u_t$，货币 ETF 的单位持仓向量为 $e_t^{\rm cash}$。辅助配置取
\[
z_t(a)=a u_t+(1-a)e_t^{\rm cash},\qquad 0\leq a\leq1.
\]
使用调仓日前的252日窗口，在共同完整观察上估计样本协方差$\widehat S_t$，并施加与现有估计模块一致的半正定数值修复。主模型权重仍由原有Ledoit--Wolf输入求得，辅助比较则对双方权重统一使用$\widehat S_t$衡量风险，选择满足
\[
z_t(a)^\mathsf T\widehat S_t z_t(a)
\leq w_t^\mathsf T\widehat S_t w_t
\]
的最大 $a$。方差计算保留货币 ETF 的波动及其与其他资产的协方差。若整个简单组合已满足上限，则保留全部简单组合；若混合集合的最低方差仍高于目标，则采用集合内最低方差点并记录未达到上限的日期。

该规则使用当期可获得的信息，保持多头、满仓和无杠杆。风险匹配针对条件估计，实际波动由随后实现收益决定，因此结果同时报告预测风险误差与实现波动。每周重新求得混合权重，随后按日漂移，并根据混合组合自身的交易量扣费。

样本矩阵尺度来自历史可行性诊断。收缩对货币估计风险的影响见第\ref{sec:lw}节；本文同时保留LW矩阵下的最低可达风险及可行性记录，两种尺度的结果分别解释。

可行性检验还发现，在主模型的Ledoit--Wolf收缩协方差下，\diagLWMatchInfeasible{}个配对决策的混合集合最低方差高于主模型目标。向共同对角尺度收缩提高了货币ETF的估计风险，使集中货币持仓在该估计下也需要承担较大方差。表中辅助结果因此采用同一历史样本协方差统一衡量双方风险，同时完整保留收缩估计的可行性记录。这一现象揭示了风险模型与工具集合的相互作用：可交易的低波动资产如何进入协方差估计，会影响可实现的风险匹配。

这里的配对记录对应\primaryRebalances{}个周度决策日与两个简单基准。样本矩阵匹配是给定历史规格的辅助比较，收缩矩阵下的不可行结果亦构成研究发现。由于匹配尺度在历史诊断中确定，所得差异应解释为该尺度下的配置证据；既不能称作所有风险维度均相同，也不能据此识别与风险模型无关的超额收益能力。

\figmaybe{primary_diagnostic_risk_matched.pdf}{主模型与事前风险匹配辅助配置的净值}{fig:diagnostic_matched}
\subsection{成本、执行与估计敏感性}
表\ref{tab:diagnostic_robustness}将稳健性检验统一限定于当前主模型。成本变化沿用相同的目标路径，执行延迟重新核算交易前权重，均值半衰期实验则保留年度参数路径。

\begin{table}[htbp]\centering
\caption{当前 Global RRP 的实施与估计敏感性}\label{tab:diagnostic_robustness}
\fitthesistable{\begin{tabular}{lrrrr}
\toprule
检验设定 & 净年化收益 & 年化波动 & 夏普 & 最大回撤\\\midrule
单边成本0 bps & 6.01\% & 3.58\% & 1.649 & -6.79\%\\
单边成本3 bps & 5.94\% & 3.58\% & 1.631 & -6.80\%\\
单边成本10 bps & 5.78\% & 3.58\% & 1.588 & -6.82\%\\
单边成本20 bps & 5.55\% & 3.58\% & 1.527 & -6.85\%\\
均值半衰期10日 & 5.86\% & 3.57\% & 1.614 & -6.83\%\\
均值半衰期40日 & 5.96\% & 3.59\% & 1.634 & -6.86\%\\
延后一交易日执行 & 5.84\% & 3.59\% & 1.600 & -6.91\%\\
\bottomrule
\end{tabular}}
\end{table}

在单边20 bps费用下，主模型净年化收益为 \diagCostTwentyReturn。将原目标顺延一个交易日执行后，净年化收益为 \diagDelayReturn，年化波动为 \diagDelayVol，夏普为 \diagDelaySharpe，最大回撤为 \diagDelayDD。费用与执行实验把交易条件变化映射为可观察的收益风险差异。10日与40日半衰期结果进一步说明，均值估计的响应速度会改变给定年度惩罚下的配置路径。

较短半衰期提高近期收益的影响，也提高了交易需求；较长半衰期使输入调整更平缓。其经济取舍是信号响应速度与估计噪声、成本之间的平衡，与DeMiguel等关于估计误差影响配置收益的讨论相联系\cite{demiguel2009optimal}。在本次敏感性范围内，净年化收益和风险指标保持相近数量级，说明主结果对所检验的均值响应速度具有一定稳定性。

\subsection{补充分析的口径与统计方法}
\begin{table}[htbp]\centering
\caption{三类分析的对象与解释口径}\label{tab:analysis_scope}
\begin{tabularx}{\textwidth}{@{}l L L@{}}\toprule
分析 & 数据与统计量 & 回答的问题\\\midrule
风险贡献 & 决策日前协方差与目标权重；波动贡献及其份额 & 给定风险模型下，各资产分担多少估计风险\\
风险匹配 & 同一历史样本矩阵衡量双方目标方差；随后比较实现绩效 & 指定事前方差尺度相近时，不同配置如何表现\\
Brinson归因 & 每日收益记账权重与实现收益；相对基准的日均基点 & 已实现收益差如何分解为配置、选择、交互和费用\\\bottomrule
\end{tabularx}
\end{table}
三类分析共享持仓路径。主模型风险贡献按LW计算，并以样本矩阵作敏感性对照；匹配统一使用样本矩阵，Brinson直接使用实现收益。方差相近不保证尾部或因子风险相同，正的归因贡献也不能替代统计推断。

\subsubsection{配对区块自助抽样}
区块抽样通过重抽连续收益片段保留局部时间依赖\cite{kunsch1989bootstrap}。本文将主模型与比较对象的日收益按同一日期配对，每次对两列使用相同的区块起点，以保留模型间同期相关性。每组进行2000次重抽样，分别采用5、21和63个交易日的区块长度；主表报告21日设定，图形展示三种长度的敏感性。

每次重抽样计算主模型减比较对象的年化夏普差，以2.5\%和97.5\%分位点构造95\%区间。夏普的收益序列与主绩效口径一致，扣除首日净值基期对应的收益观察。区间刻画固定历史规格下的配对差异不确定性，属于逐项百分位区间。Ledoit 与 Wolf 的研究讨论了学生化时间序列自助法在夏普差异推断中的作用\cite{ledoit2008sharpe}；本文采用的百分位区间用于探索性比较，结论同时考虑区块长度和历史规格选择条件。
\subsection{配对绩效差异与统计不确定性}
表\ref{tab:diagnostic_bootstrap}报告主模型减比较对象的夏普差及21日区块抽样95\%区间；正值表示主模型样本夏普较高。图\ref{fig:diagnostic_bootstrap}同时展示5日与63日区块长度的结果，用于评价对局部依赖长度的敏感性。

\begin{table}[htbp]\centering
\caption{配对夏普差异的区块抽样区间}\label{tab:diagnostic_bootstrap}
\fitthesistable{\begin{tabular}{lrr}
\toprule
比较对象 & 主模型减比较对象的夏普差 & 95\%区间\\\midrule
仅风险预算参考 & 0.412 & [0.184, 0.608]\\
收益短缺惩罚置零 & 0.067 & [-0.043, 0.171]\\
冻结评价首年参数 & -0.010 & [-0.062, 0.042]\\
Equal Weight（风险匹配） & 0.427 & [0.034, 0.805]\\
60/40 Benchmark（风险匹配） & 0.580 & [0.100, 1.034]\\
\bottomrule
\end{tabular}}
\end{table}

\figmaybe{primary_diagnostic_bootstrap.pdf}{不同区块长度下的配对夏普差及95\%区间}{fig:diagnostic_bootstrap}
\FloatBarrier

在5日、21日和63日三种区块长度下，主模型相对风险预算参考及两条风险匹配辅助路径的夏普差异区间均位于零上方。这为第二阶段优化及资产构成调整提供了方向一致的样本证据。相对收益短缺项置零和冻结首年参数，差异区间均覆盖零：前者的贡献更适合解释为收益、波动与换手之间的权衡；后者说明两种参数规则在本样本中表现接近。

本文据此将经验结论分为配置结构与参数规则两个层次。配置结构获得了消融和风险匹配的共同支持，年度参数则体现一倍标准误集合内对换手与稳定性的选择偏好。所有区间均条件于当前历史规格，属于逐项百分位区间，研究证据定位为固定资产池与既定实施规则下的探索性经验结果。
\FloatBarrier

\Needspace{9\baselineskip}
\subsection{资产配置收益归因与下跌日表现}
\label{sec:brinson}
\subsubsection{仓位、风险贡献与收益归因的分析对象}
仓位管理决定资本投向，风险贡献衡量各项持仓对组合风险的影响，收益归因则解释组合相对基准的实际收益差。本文以权重和协方差计算资产层面的波动风险贡献，并以Brinson--Fachler框架核算配置收益。Barra采用因子暴露、因子协方差和特异风险刻画组合风险，其因子收益归因还需要相应的因子收益序列\cite{msci2026factor}。因此，本文以下归因的分析单位是ETF资产组，结论对应可观察的持仓决策。

Brinson与Fachler提出了按类别分解配置和选择作用的框架\cite{brinson1985measuring}；Brinson、Hood与Beebower进一步讨论了投资政策、配置和选择决策的绩效评价\cite{brinson1986determinants}。本文采用Brinson--Fachler的类别配置口径，并依据Bacon的专业综述分别报告配置、选择与交互项\cite{bacon2019attribution}。这一框架适合回答：相对同一比较基准，多配或少配某类资产，贡献了多少实际收益差。

\subsubsection{计算口径及现金配置的含义}
设$w^P_{g,t}$与$w^B_{g,t}$分别为组合和基准在当日收益发生前的第$g$组权重，$r^P_{g,t}$与$r^B_{g,t}$为组内按资本加权的当日收益，$R^B_t=\sum_g w^B_{g,t}r^B_{g,t}$为基准毛收益。配置、组内选择与交互效应分别为
\begin{align}
A_{g,t}&=(w^P_{g,t}-w^B_{g,t})(r^B_{g,t}-R^B_t),\\
S_{g,t}&=w^B_{g,t}(r^P_{g,t}-r^B_{g,t}),\\
I_{g,t}&=(w^P_{g,t}-w^B_{g,t})(r^P_{g,t}-r^B_{g,t}).
\end{align}
满额投资使两组权重各自加总为一，故$\sum_g(A_{g,t}+S_{g,t}+I_{g,t})=R^P_t-R^B_t$。将两条路径的费用率分别记作$c^P_t$、$c^B_t$，另列费用效应$C_t=c^B_t-c^P_t$后，各项之和等于当日净收益差。

配置项衡量类别权重偏离带来的相对贡献。若组合超配货币ETF，而且货币组收益高于当日基准整体收益，则货币配置项为正。这并不要求货币ETF自身取得很高收益。例如，在一个仅含现金与股票的假设单期中，现金收益为零、股票下跌10\%，基准持有10\%现金、组合持有50\%现金，两者收益分别为$-9\%$与$-5\%$。组合领先4个百分点，其中现金配置项为3.6个百分点，股票低配项为0.4个百分点。这一例子区分了现金自身的收益贡献、现金超配的归因贡献和整个仓位调整的净效果。

组内选择项反映同一类别中ETF权重构成的差异；本文所称“选择”对应基金配置，不涉及穿透基金后的个股选择。交互项记录类别权重偏离与组内收益偏离同时出现的影响。本文将其独立列示，使三项之和可以逐日复核。该分解属于给定基准下的会计恒等式，关于配置规则作用的实验依据由前述消融与风险匹配提供。

\subsubsection{资产分组、信息时序与核算结果}
为单独观察现金、利率和信用配置，本文在归因中设置货币、国债、信用债、可转债、境内权益、境外权益与商品相关七个互斥组。该分组用于解释收益，不改变资产池或优化规则。港股计入境外权益；商品相关组沿用项目工具映射，包含煤炭等与商品相关的权益工具，因此其解释保持在ETF类别层面。

基准分别采用Equal Weight及其事前风险匹配路径。计算使用与已公布日收益一致的权重：调仓日使用当日目标，其他日期使用收益漂移后的持仓。组内收益按当日组权重归一化。若组合某组权重为零，其组收益取原始等权基准的对应组收益；若风险匹配基准某组为空，同样采用原始等权组收益作为类别参照。这一约定保持空组解释一致并保留逐日恒等式。空组的参照收益以及资产分组会影响分项数值，因此本文固定上述口径，并同时核查合计收益差。全部持有资产在归因日期均有有效收益，重建毛收益与已公布路径一致。

表\ref{tab:brinson_summary}使用日均基点表示各项贡献，一个基点等于0.01个百分点。该统计是单日效应的算术平均，适合比较全样本和条件样本；复合累计收益及年化收益仍由前文净值路径计算。Bacon讨论了多期归因需要处理收益链接的问题\cite{bacon2019attribution}，本文据此将单日均值与复合绩效分别报告。

归因样本与主绩效中由相邻净值计算的收益样本对齐，2018年1月2日作为净值基期，归因从下一交易日开始。每日四项效应之和等于当日净收益差，因而各项均值之和也等于平均日净收益差。对这些均值乘以252只会得到算术年化差异，无法直接还原两项复合年化收益之差。表中各项独立四舍五入，显示值的加总可能有末位差异，逐日未舍入数值另行核验。



全样本共\attrDays{}个交易日，沪深300ETF当日收益小于零的观察有\attrDownDays{}个。下跌日按当期实现收益事后划分，用于解释已实现损益，与表\ref{tab:position_states}中依据历史信息划分的弱势状态口径不同。相对Equal Weight，主模型在下跌日的平均净收益差为\attrEqualNetDown{}个基点，全样本对应\attrEqualNetAll{}个基点；这一组合同时体现了下跌时的损失缓冲和较低风险配置的长期收益机会成本。相对风险匹配等权路径，全样本和下跌日的平均净收益差分别为\attrMatchedNetAll{}和\attrMatchedNetDown{}个基点。

\begin{table}[htbp]\centering
\caption{各资产组的配置效应：平均单日贡献（基点）}\label{tab:brinson_allocation}
\fitthesistable{\begin{tabular}{lrrrr}
\toprule
归因组 & 等权全样本 & 等权下跌日 & 匹配全样本 & 匹配下跌日\\\midrule
货币 & -0.34 & 7.04 & 0.38 & -7.06\\
国债 & -0.58 & 14.43 & 0.04 & 4.85\\
信用债 & -0.31 & 4.75 & -0.06 & 1.42\\
可转债 & -0.00 & 0.01 & 0.01 & -0.36\\
境内权益 & 0.17 & 13.35 & 0.03 & 1.70\\
境外权益 & -0.14 & -1.08 & -0.02 & -1.02\\
商品相关 & -0.21 & -1.94 & 0.33 & -0.66\\
\bottomrule
\end{tabular}}
\end{table}
\FloatBarrier

表\ref{tab:brinson_allocation}进一步显示，相对原始等权基准，下跌日的货币、国债及信用债配置均提供正的平均配置贡献，其中货币项为\attrEqualCashDown{}个基点。风险匹配等权路径本身配置了较多货币ETF，主模型相对该路径的货币项为\attrMatchedCashDown{}个基点，国债、信用债与境内权益配置项则为正。两组结果说明，现金和债券可以通过资本配置参与下行风险管理，而具体贡献的符号取决于基准已有的持仓。主模型相对风险匹配路径的收益来源还包括组内配置和交互效应，因而需要结合完整分解评价。

两个基准承担不同解释任务：原始等权显示风险规模和配置结构共同形成的收益差，风险匹配路径显示指定事前方差尺度下的差异。下跌日分组会突出低权益敞口的损失缓冲，表中结果应与全样本及实现风险同时阅读。沪深300ETF下跌日只是境内宽基下行的观察集合，不能代表所有市场的同步危机，也不构成策略提前预测下跌的证据。
\FloatBarrier

